limiting magnitude (find?)
nickel - reliably
    time exposure as well?
    - check out limiting magnitude vs time (diminishing returns)

distance vs brightness - using hr to fit

M_limit \approx 2.5log(D^2 \times \sqrt{t}) + M 
D is diameter, M is factor based on telescope - noise

    - linear area corresponds to 5-45k counts on ccd, beyond = overexposure and bad stuff
    - probably alreayd over exposing before that
    - full saturation = 65k counts (16bit camera) 2^16
    - more color, takes longer

signal - noise
goal \frac{Signal}{Noise} = 5x (can go down to 2-3, but gets grainier)
(think of \rho and how it compares)

~40s for mag 12 on nickel telescope (TEST)
    - depends on atmosphere

!!migration in seismic (take energy from diffraction - concentrate to a point source (or similar))
area diffraction
bessel functions

focal length of secondary determines resolution (something we must find)

leakage from potential wells in the ccds (noise)
thermal electrons from other stuff (noise) - called dark current (often by the circuits heating up)

    - explore different ways to find Noise

you will get vertical noise more than anything (column connected - not row)

Signal - Noise ratio = taken by \sqrt{\delta t} - only works on linear section (stats, P(t) doesnt change)

file sharing?

explore stacks from unistellers
monte carlo? 
evscope - each pixel has a different filter - explore how stacking works?
[R G 
 G B]*n


this vs nickel, has bayer something??? - research this (LPC Electronic telescope paper)
IR is intensified - when looking at color data
    - intensities should map well, but color is where it gets weird
    - nickel for gaging distance



counting stats - how count increases with lowering thresholds (lower bound), and how that relates to Signal - Noise ratios

